\section{Immutable Parent Pointers}
\label{sec:immutable-parent-pointers}

The immutable parent pointer is the foundational structural primitive of the Recursive Spawning Protocol. It converts the delegation chain from a forensic reconstruction — assembled after the fact from mutable logs and policy assertions — into a runtime-verifiable artifact that is immutable, cryptographically enforced, and architecturally verifiable at every point in the chain's lifecycle. It renders all modes of autonomous drift structurally impossible regardless of the operational urgency, architectural complexity, or adversarial sophistication of any machine actor in the chain.

This section specifies: the formal definition of the parent pointer relation and its four structural constraints; the chain operator $\Chain{\cdot}$ and its recursive definition with the human anchor base case; the Fact Layer write protocol that enforces immutability at the protocol level; the Purpose Layer's structurally necessary role in spawn authorization; the chain traversal algorithm with correctness arguments; and a summary framing immutability as an architectural constraint rather than a policy choice.

% ------------------------------------------------------------------
\subsection{Formal Definition: ParentPointer}
\label{sec:parent-pointer-formal}

\begin{dfn}[ParentPointer Relation]
\label{dfn:parent-pointer}
Let $\mathcal{N}$ be the set of all agent nodes in a \bxthree{} system implementing the Recursive Spawning Protocol. For any node $c \in \mathcal{N}$ spawned via a valid RSP spawn event, the \emph{parent pointer} $\ParentPointer{p}{c}$ is the immutable, cryptographically signed reference from child $c$ to its immediate parent $p \in \mathcal{N}$, established exactly once at node provisioning via a Fact Layer atomic write operation. The parent pointer relation satisfies four structural constraints:

\begin{enumerate}
    \item[(P1) \textbf{Uniqueness}] For every node $c \in \mathcal{N}$ spawned under RSP, there exists exactly one node $p \in \mathcal{N}$ such that $\ParentPointer{p}{c}$ holds at all times after provisioning. No node may have multiple simultaneous parent pointers. No parent pointer may be removed, redirected, or reassigned after establishment.

    \item[(P2) \textbf{Immutability}] After the provisioning write confirming $\ParentPointer{p}{c}$ is recorded in the Fact Layer ledger, no operation exists — from any source, with any privilege level, under any system condition — that can modify or delete $\ParentPointer{p}{c}$. The pointer is permanent for the operational lifetime of $c$. This is not a policy enforced by access control; it is a structural property of the protocol. The modification operation is not defined.

    \item[(P3) \textbf{Directionality}] The parent pointer is an intrinsic property of the child node, not a reference held by the parent node. The child $c$ carries $\ParentPointer{p}{c}$ as an immutable attribute in its own Fact Layer state record. The parent $p$ holds no mutable reference to $c$. This means the chain is traversable from any node outward to its complete lineage regardless of the current operational state of any other node in the chain. A parent that has terminated, crashed, or been decommissioned does not affect $c$'s ability to reference its parent through $\ParentPointer{p}{c}$.

    \item[(P4) \textbf{Human Anchor Termination}] The root of every RSP chain is a human principal node $h$ for which $\ParentPointer{\bot}{h}$ holds — the null parent marker indicating that $h$ is the ultimate delegation origin. No machine actor in the RSP model may serve as a chain root. The root human anchor is designated by the Purpose Layer at system initialization and recorded in the Fact Layer authorization ledger.
\end{enumerate}
\end{dfn}

\medskip
\noindent\textbf{Notation.} We write $\ParentPointer{p}{c}$ to denote the parent pointer relation between parent $p$ and child $c$. The equivalent functional notation is $\alpha(c) = p$. The notation $\alpha(c) = \bot$ denotes that node $c$ has no parent — it is the root human anchor, the ultimate delegation origin. The functions $p = \parent(c)$ and $c \in \children(p)$ are used interchangeably with the pointer relation. The notation $\alpha^{(j)}(c)$ denotes the $j$-fold application of $\alpha$, and $\alpha^{(0)}(c) = c$.

\begin{dfn}[Spawn Event]
\label{dfn:spawn-event}
A \emph{spawn event} is the 5-tuple
\begin{equation}
\mathrm{spawn}(c,\ p,\ \sigma_c,\ t,\ \phi)
\end{equation}
where $c$ is the newly created child node, $p \in \mathcal{N}$ is the authorizing parent node, $\sigma_c$ is the authorized execution scope, $t$ is the wall-clock timestamp, and $\phi$ is the cryptographic signature produced by the Purpose Layer over the spawn request. The Fact Layer atomically: (1) creates node $c$, (2) sets $\ParentPointer{p}{c}$, (3) sets the node's authorized scope to $\sigma_c$, (4) seals the memory envelope, and (5) appends the event to the forensic ledger $\mathcal{L}$.
\end{dfn}

The initialization of $\ParentPointer{p}{c}$ occurs atomically during the spawn event. The Fact Layer records the spawn event in the forensic ledger and writes $\ParentPointer{p}{c}$ into the node's sealed memory envelope as part of the same atomic transaction. This separation of authorization (Purpose Layer) from assignment (Fact Layer) is essential: it prevents a parent from claiming to have spawned a node whose parent pointer was set to a different actor, and it prevents a child from falsifying its own ancestry after provisioning.

\begin{prop}[Immutability of ParentPointer]
\label{prop:parent-pointer-immutable}
For any child node $c$ with parent $p$, once the spawn event $\mathrm{spawn}(c, p, \sigma_c, t, \phi)$ has been recorded in the forensic ledger $\mathcal{L}$, the value $\ParentPointer{p}{c}$ is fixed for the operational lifetime of $c$. No subsequent operation — including operations issued by $p$, by any Purpose Layer actor, by the Bounds Engine, or by $c$ itself — can alter $\ParentPointer{p}{c}$.
\end{prop}

This property distinguishes ParentPointer from conventional mutable parent references in multi-agent systems. A mutable reference can be overwritten by a subsequent administrative operation; an immutable pointer cannot. The distinction is not behavioral — it is architectural.

% ------------------------------------------------------------------
\subsection{The Chain Operator and Human Anchor Base Case}
\label{sec:chain-operator-definition}

The chain operator $\Chain{\cdot}$ is the primary tool for chain traversal, verification, and accountability attribution. It constructs the complete delegation lineage from any node to the human anchor by recursively applying the parent pointer function.

\begin{dfn}[Chain Operator]
\label{def:chain-operator}
For any agent node $a \in \mathcal{N}$ in a Recursive Spawning chain, the \emph{chain} $\Chain{a}$ is defined recursively as:

\begin{equation}
\Chain{a} =
\begin{cases}
\ParentPointer{\parent(a)}{a} \ \cup \ \Chain{\parent(a)} & \text{if } \alpha(a) \neq \bot \\[10pt]
\{\, a \,\} & \text{if } \alpha(a) = \bot \quad \text{(root human anchor)}
\end{cases}
\label{eq:chain-recursive}
\end{equation}

Equivalently, expanding the recursion to reveal the full lineage:

\begin{equation}
\Chain{a} = \bigl\{\, a,\ \alpha(a),\ \alpha^{(2)}(a),\ \alpha^{(3)}(a),\ \ldots,\ \alpha^{(k)}(a) \,\bigr\}
\label{eq:chain-expanded}
\end{equation}

where $k$ is the depth of node $a$ in the chain (the number of successive parent applications required to reach the root), $\alpha^{(j)}(a)$ denotes the $j$-fold application of $\alpha$, and $\alpha^{(k)}(a)$ is the unique root node $n_0$ satisfying $\alpha(n_0) = \bot$ and $\hmannchor{n_0}$.
\end{dfn}

The chain $\Chain{a}$ is the ordered set of all nodes on the delegation path from $a$ to and including the human anchor. It includes the terminal human anchor node, which is the only node in any RSP chain with $\alpha(n) = \bot$. The chain is always finite: by the RSP termination property, every RSP chain has depth at most $D_{\max}$, so the recursion in Equation~\ref{eq:chain-recursive} terminates at the root within at most $D_{\max} + 1$ successive applications of $\alpha$.

\medskip
\noindent\textbf{Base case: root human anchor.} The root of every RSP chain is a Purpose Layer entity designated as a human principal, established at system initialization when a human principal $h$ authorizes the first agent node $n_0$. The Purpose Layer records $h(n_0) = h$ and $\alpha(n_0) = \bot$ in the Fact Layer authorization ledger. This is the only valid chain-root in RSP: no machine actor may serve as a chain origin, and no node may have $\alpha(n) = \bot$ unless it carries the $\hmannchor{\cdot}$ designation. The non-collapsing human anchor property ensures $h(n) = h(n_0)$ for all nodes in all chains spawned from $n_0$.

The chain operator satisfies three critical properties by construction:

\begin{enumerate}
\item[(C1) \textbf{Uniqueness}] For any node $a$, $\Chain{a}$ is a singleton set constructed by iterating $\alpha$ from $a$ to the root. There is exactly one such set; no alternative chain exists for any node in a correctly provisioned RSP system.

\item[(C2) \textbf{Determinism}] $\Chain{a}$ is identical at any runtime point to what it was at the moment of $a$'s provisioning. Because $\alpha(c)$ is immutable for all $c$, no retroactive chain modification is possible. The chain is not reconstructed from logs — it is directly computable by iterating $\alpha$ from any node until reaching $\bot$.

\item[(C3) \textbf{Human termination}] Every $\Chain{a}$ terminates at a human principal. The RSP protocol does not permit machine actors to serve as chain roots. This property is structurally enforced by the Fact Layer pre-commit hook, which rejects any spawn event whose root is not a registered human principal.
\end{enumerate}

% ------------------------------------------------------------------
\subsection{Immutability Enforcement: Fact Layer Write Protocol}
\label{sec:fact-layer-write-protocol}

The immutability of $\ParentPointer{p}{c}$ is not enforced by access control — it is a structural property of the protocol established by the Fact Layer write protocol. The Fact Layer is the accountability infrastructure layer in the \bxthree{} architecture: it maintains the authoritative record of all authorization state, including parent pointers, warrants, and spawn events. The write protocol defines the valid state transitions for Fact Layer records and explicitly excludes the modification of $\alpha(n)$ after provisioning.

The Fact Layer write protocol enforces immutability through four mechanisms:

\begin{enumerate}
    \item[(W1) \textbf{No modification operation defined.}] The Fact Layer write protocol defines exactly two state transitions for $\ParentPointer{p}{c}$: its initial creation during the atomic spawn event, and its permanent retention thereafter. No state transition for modifying or deleting $\ParentPointer{p}{c}$ is defined in the protocol specification. An operation that is not defined cannot be executed; it cannot be authorized, bypassed, or emulated.

    \item[(W2) \textbf{Pre-commit hook rejection.}] The Fact Layer pre-commit hook intercepts any provisioning request that attempts to overwrite an existing parent pointer. The hook inspects the Fact Layer state record for $c$; if $\ParentPointer{p}{c}$ is already established and the incoming request attempts to set $\ParentPointer{p'}{c}$ where $p' \neq p$, the hook rejects the request before any write occurs. The rejection is unconditional: the hook does not evaluate credentials, privilege levels, or operational urgency.

    \item[(W3) \textbf{Sealed memory envelope.}] The parent pointer field $\ParentPointer{p}{c}$ is encrypted in the node's persistent memory envelope using a Fact Layer-only symmetric key. Even a compromised Bounds Engine cannot read $\ParentPointer{p}{c}$ because it never has access to the decryption key. The envelope is HMAC-signed; any modification to the ciphertext invalidates the signature. The Fact Layer decrypts the envelope only when reading $\ParentPointer{p}{c}$ to service a chain-traversal or audit request, and re-encrypts and re-signs after each read, preventing replay attacks based on stale plaintext.

    \item[(W4) \textbf{Atomic provisioning.}] The parent pointer $\ParentPointer{p}{c}$ and its Purpose Layer warrant $\phi$ are created as a single atomic Fact Layer write. There is no temporal window in which the pointer exists without authorization, or in which authorization exists without a corresponding pointer.
\end{enumerate}

\medskip
\noindent\textbf{Why access control is insufficient.} The distinction between access-control-based immutability and protocol-level immutability is the distinction between a promise and a guarantee. In an access-control model, immutability is enforced by restricting which principals may issue modification operations. If an attacker acquires sufficient privileges — through credential theft, insider compromise, or software vulnerability — the access control barrier is circumvented and immutability is violated. The immutability was a policy, enforceable until it was not.

In the Fact Layer write protocol, the immutability of $\ParentPointer{p}{c}$ is a property of the protocol itself. There is no operation defined for modifying $\ParentPointer{p}{c}$ after provisioning, so no amount of privilege elevation, credential compromise, or software vulnerability can produce a valid modification operation. The protocol does not enforce immutability — it makes immutability the only valid state transition. The distinction is architectural, not a matter of degree.

This is the precise failure mode that the Fact Layer write protocol eliminates: retroactive manipulation of chain topology for the purpose of accountability laundering. Without protocol-level immutability, an adversarial actor could re-parent a node to a favorable parent after the fact, obscuring the original authorization chain and making drift detection impossible. With protocol-level immutability, this manipulation is structurally impossible.

% ------------------------------------------------------------------
\subsection{Spawn Authorization: Purpose Layer Role}
\label{sec:purpose-layer-authorization}

The Purpose Layer plays two structurally necessary roles in the parent pointer lifecycle: it originates the spawn authorization policy that defines the conditions under which a new parent pointer may be created, and it issues the spawn warrant that is the prerequisite for the Fact Layer provisioning write. These two roles are inseparable from the parent pointer's immutability: if the Purpose Layer could issue warrants after the fact, or if a machine actor could create a parent pointer without a Purpose Layer warrant, the chain's integrity guarantees would be void.

\medskip
\noindent\textbf{Authorization policy origination.} At system initialization, the Purpose Layer defines the spawn authorization policy $P_{\mathrm{spawn}}$ and records it in the Fact Layer authorization ledger. $P_{\mathrm{spawn}}$ specifies the mandatory conditions for any valid spawn event:

\begin{enumerate}
    \item[(S1) \textbf{Scope refinement.}] The child scope must be a strict subset of the parent scope: $R(n_{\mathrm{child}}) \subset R(n_{\mathrm{parent}})$. No lateral expansion, no superset, no equality. This ensures that each spawn event narrows the delegation scope, preserving the intent-refinement property of the chain.

    \item[(S2) \textbf{Operational necessity.}] The parent must declare, in the Fact Layer authorization ledger, the precise condition requiring child spawning. The declaration must identify why the condition cannot be resolved within $R(n_{\mathrm{parent}})$ and why the child is the minimal delegation action required to address it.

    \item[(S3) \textbf{Human anchor consistency.}] The spawn event must not alter the human anchor $h(n)$ of any node in the chain. The child inherits $h(n_{\mathrm{child}}) = h(n_{\mathrm{parent}}) = h(n_0)$; the anchor is non-collapsing across all spawn events.
\end{enumerate}

\medskip
\noindent\textbf{Warrant issuance as exclusive origin.} The Purpose Layer warrant $\phi$ is the sole authorized precondition for the Fact Layer provisioning write that establishes $\ParentPointer{p}{c}$. No machine actor may initiate a spawn event without a matching warrant in the Fact Layer ledger — the Fact Layer pre-commit hook rejects any provisioning request that lacks a valid warrant before the operation reaches the ledger write stage. This makes the Purpose Layer's exclusive authorization origin role a structural guarantee, not a configuration option.

The warrant is cryptographically signed by the Purpose Layer and carries the spawn parameters: child node identity, authorized scope $\sigma_c$, parent identity, and timestamp. The Fact Layer pre-commit hook verifies the signature against the registered Purpose Layer public key before accepting the provisioning write. This prevents unauthorized chain construction even if all other RSP mechanisms fail simultaneously.

\begin{prop}[Exclusive Authorization Origin]
\label{prop:exclusive-auth-origin}
The Purpose Layer is the sole origin of spawn authorization in RSP. A machine actor cannot self-authorize a spawn event, cannot retroactively authorize an unauthorized spawn, and cannot modify the chain topology after provisioning. Any attempt to establish $\ParentPointer{p}{c}$ without a matching Purpose Layer warrant $\phi$ is rejected at the Fact Layer pre-commit hook, before the operation reaches the ledger.

The Purpose Layer's exclusive role as authorization origin prevents the chain from being bootstrapped by a machine actor acting autonomously. Even if every other RSP mechanism failed simultaneously, the sole-authorization origin property would prevent a machine actor from independently constructing a delegation chain — the ghost authority failure mode is structurally blocked.
\end{prop}

% ------------------------------------------------------------------
\subsection{Chain Traversal Algorithm}
\label{sec:chain-traversal}

The depth bound $D_{\max}$ provides a constant-time upper bound on chain traversal cost. The worst-case number of pointer dereferences required to reach the human anchor from any node is $D_{\max} + 1$, independent of the total number of nodes in the system. A system with 10,000 spawned nodes and $D_{\max} = 5$ requires at most six pointer dereferences to reach a human — the same as a system with 10 nodes. Mutable parent pointers would allow the chain to be extended retroactively, making the effective depth potentially unbounded and the worst-case escalation time unbounded as well. Immutability is what makes the depth bound a structural guarantee, not merely a policy.

\begin{algorithm}
\caption{Chain-Traverse($a$) — Retrieve the delegation chain from node $a$ to human anchor}
\label{alg:chain-traverse}
\begin{algorithmic}[1]
\REQUIRE $a$ is a provisioned RSP node
\ENSURE Ordered list $\Chain{a} = \langle n_0, n_1, \ldots, n_k \rangle$ with $n_0 = a$ and $n_k = h(a)$

\STATE $chain \leftarrow \text{empty list}$
\STATE $current \leftarrow a$ \COMMENT{Start at the query node}

\WHILE{$\text{true}$}
    \STATE $\text{append}(chain,\ current)$ \COMMENT{Add current node to chain}
    \IF{$\alpha(\text{current}) = \bot$}
        \STATE \textbf{break} \COMMENT{Root human anchor reached — chain complete}
    \ENDIF
    \STATE $\text{current} \leftarrow \alpha(\text{current})$ \COMMENT{Move to parent}
\ENDWHILE

\RETURN $chain$
\end{algorithmic}
\end{algorithm}

\medskip
\noindent\textbf{Correctness arguments.}

\textit{Termination.} The algorithm terminates because RSP chains are depth-bounded by construction. Each iteration of the while-loop moves $\text{current}$ to the parent node $\alpha(\text{current})$, which is strictly closer to the root along the chain. Since the chain has finite depth at most $D_{\max} + 1$, the loop executes at most $D_{\max} + 1$ iterations before $\alpha(\text{current}) = \bot$ is reached. The loop therefore terminates in finite time. This argument holds regardless of the current operational state of any node in the chain, including nodes that may have terminated or entered error states, because the traversal only reads the immutable $\alpha$ values.

\textit{Correctness of the returned set.} By Definition~\ref{def:chain-operator}, $\Chain{a}$ is the set containing $a$, $\alpha(a)$, $\alpha(\alpha(a))$, and so on, until the root. The algorithm produces exactly this sequence: it starts at $a$, appends each node, moves to its parent via $\alpha$, and stops when $\alpha(\text{current}) = \bot$ is reached. The output is identical to $\Chain{a}$ by construction.

\textit{Uniqueness of result.} The algorithm produces a unique result for any input node because $\alpha$ is a total function (defined for every provisioned node) and is injective on the chain direction (every node has exactly one parent, and the root has no parent). There is no branching, no alternative path, and no non-deterministic choice at any step in the algorithm.

\textit{Immutability guarantee.} The algorithm reads $\alpha$ values but never writes them. The traversal is read-only and does not interact with the Fact Layer write protocol. The algorithm can be executed at any runtime point without affecting the immutability of the chain it traverses.

\begin{algorithm}
\caption{Chain-Verify($n$) — Verify the authorization chain from node $n$ to human anchor}
\label{alg:chain-verify}
\begin{algorithmic}[1]
\REQUIRE $n$ is a provisioned RSP node
\ENSURE $\text{Boolean}$ indicating whether the chain is a structurally valid RSP chain

\STATE $chain \leftarrow \text{Chain-Traverse}(n)$
\STATE $m \leftarrow \text{len}(chain)$

\STATE \COMMENT{\textit{Step V1: Verify chain terminates at human anchor}}
\IF{$\alpha(chain[m-1]) \neq \bot$}
    \STATE \RETURN $\text{false}$ \COMMENT{Root is not $\bot$ — malformed chain}
\ENDIF
\IF{$\neg\hmannchor{chain[m-1]}$}
    \STATE \RETURN $\text{false}$ \COMMENT{Root is not a designated human principal}
\ENDIF

\STATE \COMMENT{\textit{Step V2: Verify each spawn event has a Purpose Layer warrant}}
\FOR{$i \leftarrow 0$ \TO $m-2$}
    \STATE $child \leftarrow chain[i]$
    \STATE $parent \leftarrow chain[i+1]$
    \IF{$\neg\text{warrant\_exists}(child,\ parent)$}
        \STATE \RETURN $\text{false}$ \COMMENT{Missing Purpose Layer warrant for spawn}
    \ENDIF
\ENDFOR

\STATE \COMMENT{\textit{Step V3: Verify scope refinement at each link}}
\FOR{$i \leftarrow 0$ \TO $m-2$}
    \STATE $child \leftarrow chain[i]$
    \STATE $parent \leftarrow chain[i+1]$
    \IF{$\neg(R(child) \subset R(parent))$}
        \STATE \RETURN $\text{false}$ \COMMENT{Scope refinement constraint violated}
    \ENDIF
\ENDFOR

\STATE \COMMENT{\textit{Step V4: Verify depth bound compliance at each node}}
\FOR{$i \leftarrow 0$ \TO $m-1$}
    \IF{$\text{depth}(chain[i]) > D_{\max}$}
        \STATE \RETURN $\text{false}$ \COMMENT{Depth bound exceeded at node}
    \ENDIF
\ENDFOR

\RETURN $\text{true}$
\end{algorithmic}
\end{algorithm}

Chain-Verify returns \texttt{true} if and only if all four verification conditions hold simultaneously. The verification inspects structural chain properties rather than current operational state. Each condition corresponds to a fundamental RSP constraint: V1 confirms the chain terminates at a human anchor (not at a machine actor or orphaned node), V2 confirms every spawn was authorized by the Purpose Layer (not self-bootstrapped by a machine actor), V3 confirms scope refinement was respected at every link (not lateral expansion), and V4 confirms the depth bound was respected throughout (not approaching infinity). A chain that passes Chain-Verify is, by the RSP definitions, a fully authorized delegation chain terminating at a human principal.

The algorithm is deterministic and produces the same result regardless of which node in the chain initiates verification. There is no self-serving interpretation of authorization that a drifted or orphaned node could produce, because Chain-Verify inspects the complete chain topology and warrants in the Fact Layer ledger, not the node's own claims about its authorization. Chain verification is a property of the chain topology, not a property of the node being verified.

\textit{Computational cost.} Chain-Traverse runs in $O(d)$ time where $d$ is the chain depth (at most $D_{\max} + 1$). Chain-Verify runs in $O(m)$ time for the traversal plus $O(m)$ for each of the three verification loops, yielding $O(m)$ total where $m$ is chain length (also bounded by $D_{\max} + 1$). The algorithms are linear in chain length and do not depend on the total number of nodes in the system.

% ------------------------------------------------------------------
\subsection{Summary: Immutability as Architectural Constraint}
\label{sec:immutability-summary}

The immutable parent pointer is the load-bearing constraint of the Recursive Spawning Protocol. Without it, the chain is mutable — subject to retroactive modification by actors with sufficient privilege — and the RSP's correctness properties (drift prevention, chain integrity, termination) collapse to policy conventions rather than structural guarantees.

The immutability is architectural — not a stronger access control policy — because it is enforced by the Fact Layer write protocol, which defines no valid modification operation for $\ParentPointer{p}{c}$ after provisioning. This is distinct from access-control-based immutability in four ways: there is no override path for any actor; the immutability holds under all system failure conditions including those that disable other enforcement mechanisms; all sources receive equivalent enforcement treatment regardless of privilege level; and the parent pointer and its Purpose Layer warrant are created atomically with no temporal window for inconsistency.

The chain operator $\Chain{a}$ is uniquely determined by the immutable parent pointers in the chain. The chain traversal algorithm computes this operator correctly and terminates. Chain-Verify provides the runtime verification procedure for chain structural validity. Together, these mechanisms make the delegation chain an immutable, verifiable, and auditable runtime artifact — the foundation on which all other RSP correctness properties are built.